\noindent\emph{Hecke normalizations.}
For the point-level geometric quotient correspondence the congruence relation is
\[
 T_p^{\mathrm{pt}}=F_p+p\langle p\rangle F_p^{-1}.
\]
Thus $T_p^{\mathrm{pt}}-1-p\langle p\rangle$ kills rational torsion.
For the multiplicative-level geometric quotient correspondence it is
\[
 T_p^{\mu}=\langle p\rangle F_p+pF_p^{-1},
\]
with annihilator $T_p^\mu-\langle p\rangle-p$.
The quotient by the universal point and Cartier duality identify the point
and multiplicative models over $\mathbb Q$, reversing the diamond.
The canonical $q$-expansions use the latter model. The global unit
identities verify that its diamond matrix for $5$ is point multiplication
by $5^{-1}$.

This distinction is visible at a rational point-level cusp represented by
$E(q^{57}),P=q$. Its image on the multiplicative model has $j$-width $1$.
The actual point-level $5$-isogenies give $T_5c=c+5\langle5\rangle c$.
In the genus-$9$ quotient, the difference of the two such rational cusps
has order $14$. Therefore the two annihilator formulas cannot be
interchanged on the same geometric correspondence.

The clean proof uses the point-level divisor coefficients $[1,4,-5]$.
Both arrangements $[1,4,-5]$ and $[5,-4,-1]$ have been independently tested
and are nonprincipal in all three retained cases. The latter tests remain
valid numerical checks but are not used as the point-level necessity step.
The binary formal-immersion operators are unchanged, since the two odd-prime
formulas have the same reduction modulo $2$.
